\IfFileExists{stacks-project.cls}{%
\documentclass{stacks-project}
}{%
\documentclass{amsart}
}

% For dealing with references we use the comment environment
\usepackage{verbatim}
\newenvironment{reference}{\comment}{\endcomment}
%\newenvironment{reference}{}{}
\newenvironment{slogan}{\comment}{\endcomment}
\newenvironment{history}{\comment}{\endcomment}

% For commutative diagrams we use Xy-pic
\usepackage[all]{xy}

% We use 2cell for 2-commutative diagrams.
\xyoption{2cell}
\UseAllTwocells

% We use multicol for the list of chapters between chapters
\usepackage{multicol}

% This is generally recommended for better output
\usepackage{lmodern}
\usepackage[T1]{fontenc}

% For cross-file-references
\usepackage{xr-hyper}

% Package for hypertext links:
\usepackage{hyperref}

% For any local file, say "hello.tex" you want to link to please
% use \externaldocument[hello-]{hello}
\externaldocument[introduction-]{introduction}
\externaldocument[conventions-]{conventions}
\externaldocument[sets-]{sets}
\externaldocument[categories-]{categories}
\externaldocument[topology-]{topology}
\externaldocument[sheaves-]{sheaves}
\externaldocument[sites-]{sites}
\externaldocument[stacks-]{stacks}
\externaldocument[fields-]{fields}
\externaldocument[algebra-]{algebra}
\externaldocument[brauer-]{brauer}
\externaldocument[homology-]{homology}
\externaldocument[derived-]{derived}
\externaldocument[simplicial-]{simplicial}
\externaldocument[more-algebra-]{more-algebra}
\externaldocument[smoothing-]{smoothing}
\externaldocument[modules-]{modules}
\externaldocument[sites-modules-]{sites-modules}
\externaldocument[injectives-]{injectives}
\externaldocument[cohomology-]{cohomology}
\externaldocument[sites-cohomology-]{sites-cohomology}
\externaldocument[dga-]{dga}
\externaldocument[dpa-]{dpa}
\externaldocument[sdga-]{sdga}
\externaldocument[hypercovering-]{hypercovering}
\externaldocument[schemes-]{schemes}
\externaldocument[constructions-]{constructions}
\externaldocument[properties-]{properties}
\externaldocument[morphisms-]{morphisms}
\externaldocument[coherent-]{coherent}
\externaldocument[divisors-]{divisors}
\externaldocument[limits-]{limits}
\externaldocument[varieties-]{varieties}
\externaldocument[topologies-]{topologies}
\externaldocument[descent-]{descent}
\externaldocument[perfect-]{perfect}
\externaldocument[more-morphisms-]{more-morphisms}
\externaldocument[flat-]{flat}
\externaldocument[groupoids-]{groupoids}
\externaldocument[more-groupoids-]{more-groupoids}
\externaldocument[etale-]{etale}
\externaldocument[chow-]{chow}
\externaldocument[intersection-]{intersection}
\externaldocument[pic-]{pic}
\externaldocument[weil-]{weil}
\externaldocument[adequate-]{adequate}
\externaldocument[dualizing-]{dualizing}
\externaldocument[duality-]{duality}
\externaldocument[discriminant-]{discriminant}
\externaldocument[derham-]{derham}
\externaldocument[local-cohomology-]{local-cohomology}
\externaldocument[algebraization-]{algebraization}
\externaldocument[curves-]{curves}
\externaldocument[resolve-]{resolve}
\externaldocument[models-]{models}
\externaldocument[functors-]{functors}
\externaldocument[equiv-]{equiv}
\externaldocument[pione-]{pione}
\externaldocument[etale-cohomology-]{etale-cohomology}
\externaldocument[proetale-]{proetale}
\externaldocument[relative-cycles-]{relative-cycles}
\externaldocument[more-etale-]{more-etale}
\externaldocument[trace-]{trace}
\externaldocument[crystalline-]{crystalline}
\externaldocument[spaces-]{spaces}
\externaldocument[spaces-properties-]{spaces-properties}
\externaldocument[spaces-morphisms-]{spaces-morphisms}
\externaldocument[decent-spaces-]{decent-spaces}
\externaldocument[spaces-cohomology-]{spaces-cohomology}
\externaldocument[spaces-limits-]{spaces-limits}
\externaldocument[spaces-divisors-]{spaces-divisors}
\externaldocument[spaces-over-fields-]{spaces-over-fields}
\externaldocument[spaces-topologies-]{spaces-topologies}
\externaldocument[spaces-descent-]{spaces-descent}
\externaldocument[spaces-perfect-]{spaces-perfect}
\externaldocument[spaces-more-morphisms-]{spaces-more-morphisms}
\externaldocument[spaces-flat-]{spaces-flat}
\externaldocument[spaces-groupoids-]{spaces-groupoids}
\externaldocument[spaces-more-groupoids-]{spaces-more-groupoids}
\externaldocument[bootstrap-]{bootstrap}
\externaldocument[spaces-pushouts-]{spaces-pushouts}
\externaldocument[spaces-chow-]{spaces-chow}
\externaldocument[groupoids-quotients-]{groupoids-quotients}
\externaldocument[spaces-more-cohomology-]{spaces-more-cohomology}
\externaldocument[spaces-simplicial-]{spaces-simplicial}
\externaldocument[spaces-duality-]{spaces-duality}
\externaldocument[formal-spaces-]{formal-spaces}
\externaldocument[restricted-]{restricted}
\externaldocument[spaces-resolve-]{spaces-resolve}
\externaldocument[formal-defos-]{formal-defos}
\externaldocument[defos-]{defos}
\externaldocument[cotangent-]{cotangent}
\externaldocument[examples-defos-]{examples-defos}
\externaldocument[algebraic-]{algebraic}
\externaldocument[examples-stacks-]{examples-stacks}
\externaldocument[stacks-sheaves-]{stacks-sheaves}
\externaldocument[criteria-]{criteria}
\externaldocument[artin-]{artin}
\externaldocument[quot-]{quot}
\externaldocument[stacks-properties-]{stacks-properties}
\externaldocument[stacks-morphisms-]{stacks-morphisms}
\externaldocument[stacks-limits-]{stacks-limits}
\externaldocument[stacks-cohomology-]{stacks-cohomology}
\externaldocument[stacks-perfect-]{stacks-perfect}
\externaldocument[stacks-introduction-]{stacks-introduction}
\externaldocument[stacks-more-morphisms-]{stacks-more-morphisms}
\externaldocument[stacks-geometry-]{stacks-geometry}
\externaldocument[moduli-]{moduli}
\externaldocument[moduli-curves-]{moduli-curves}
\externaldocument[examples-]{examples}
\externaldocument[exercises-]{exercises}
\externaldocument[guide-]{guide}
\externaldocument[desirables-]{desirables}
\externaldocument[coding-]{coding}
\externaldocument[obsolete-]{obsolete}
\externaldocument[fdl-]{fdl}
\externaldocument[index-]{index}

% Theorem environments.
%
\theoremstyle{plain}
\newtheorem{theorem}[subsection]{Theorem}
\newtheorem{proposition}[subsection]{Proposition}
\newtheorem{lemma}[subsection]{Lemma}

\theoremstyle{definition}
\newtheorem{definition}[subsection]{Definition}
\newtheorem{example}[subsection]{Example}
\newtheorem{exercise}[subsection]{Exercise}
\newtheorem{situation}[subsection]{Situation}

\theoremstyle{remark}
\newtheorem{remark}[subsection]{Remark}
\newtheorem{remarks}[subsection]{Remarks}

\numberwithin{equation}{subsection}

% Macros
%
\def\lim{\mathop{\mathrm{lim}}\nolimits}
\def\colim{\mathop{\mathrm{colim}}\nolimits}
\def\Spec{\mathop{\mathrm{Spec}}}
\def\Hom{\mathop{\mathrm{Hom}}\nolimits}
\def\Ext{\mathop{\mathrm{Ext}}\nolimits}
\def\SheafHom{\mathop{\mathcal{H}\!\mathit{om}}\nolimits}
\def\SheafExt{\mathop{\mathcal{E}\!\mathit{xt}}\nolimits}
\def\Sch{\mathit{Sch}}
\def\Mor{\mathop{\mathrm{Mor}}\nolimits}
\def\Ob{\mathop{\mathrm{Ob}}\nolimits}
\def\Sh{\mathop{\mathit{Sh}}\nolimits}
\def\NL{\mathop{N\!L}\nolimits}
\def\CH{\mathop{\mathrm{CH}}\nolimits}
\def\proetale{{pro\text{-}\acute{e}tale}}
\def\etale{{\acute{e}tale}}
\def\QCoh{\mathit{QCoh}}
\def\Ker{\mathop{\mathrm{Ker}}}
\def\Im{\mathop{\mathrm{Im}}}
\def\Coker{\mathop{\mathrm{Coker}}}
\def\Coim{\mathop{\mathrm{Coim}}}

% Boxtimes
%
\DeclareMathSymbol{\boxtimes}{\mathbin}{AMSa}{"02}

%
% Macros for moduli stacks/spaces
%
\def\QCohstack{\mathcal{QC}\!\mathit{oh}}
\def\Cohstack{\mathcal{C}\!\mathit{oh}}
\def\Spacesstack{\mathcal{S}\!\mathit{paces}}
\def\Quotfunctor{\mathrm{Quot}}
\def\Hilbfunctor{\mathrm{Hilb}}
\def\Curvesstack{\mathcal{C}\!\mathit{urves}}
\def\Polarizedstack{\mathcal{P}\!\mathit{olarized}}
\def\Complexesstack{\mathcal{C}\!\mathit{omplexes}}
% \Pic is the operator that assigns to X its picard group, usage \Pic(X)
% \Picardstack_{X/B} denotes the Picard stack of X over B
% \Picardfunctor_{X/B} denotes the Picard functor of X over B
\def\Pic{\mathop{\mathrm{Pic}}\nolimits}
\def\Picardstack{\mathcal{P}\!\mathit{ic}}
\def\Picardfunctor{\mathrm{Pic}}
\def\Deformationcategory{\mathcal{D}\!\mathit{ef}}
\usepackage{cleveref}
\crefname{originaldefinition}{original definition}{original definitions}
\Crefname{originaldefinition}{Original definition}{Original definitions}
\makeatletter
\newlabel{situation-coherent}{{08KB}{source}{}{sourcepointer-1}{}}
\newlabel{situation-coherent@cref}{{[originaldefinition][0][]08KB}{[1][0][]source}{}{}{}}
\makeatother
\crefname{originallemma}{original lemma}{original lemmas}
\Crefname{originallemma}{Original lemma}{Original lemmas}
\makeatletter
\newlabel{more-morphisms-lemma-pushout-along-thickening}{{07RT}{source}{}{sourcepointer-2}{}}
\newlabel{more-morphisms-lemma-pushout-along-thickening@cref}{{[originallemma][0][]07RT}{[1][0][]source}{}{}{}}
\makeatother
\crefname{originallemma}{original lemma}{original lemmas}
\Crefname{originallemma}{Original lemma}{Original lemmas}
\makeatletter
\newlabel{more-morphisms-lemma-equivalence-categories-schemes-over-pushout-flat}{{07RX}{source}{}{sourcepointer-3}{}}
\newlabel{more-morphisms-lemma-equivalence-categories-schemes-over-pushout-flat@cref}{{[originallemma][0][]07RX}{[1][0][]source}{}{}{}}
\makeatother
\crefname{originallemma}{original lemma}{original lemmas}
\Crefname{originallemma}{Original lemma}{Original lemmas}
\makeatletter
\newlabel{more-morphisms-lemma-scheme-over-pushout-flat-modules}{{08KU}{source}{}{sourcepointer-4}{}}
\newlabel{more-morphisms-lemma-scheme-over-pushout-flat-modules@cref}{{[originallemma][0][]08KU}{[1][0][]source}{}{}{}}
\makeatother
\crefname{originallemma}{original lemma}{original lemmas}
\Crefname{originallemma}{Original lemma}{Original lemmas}
\makeatletter
\newlabel{more-morphisms-lemma-flat-and-free-at-point-fibre}{{080Q}{source}{}{sourcepointer-5}{}}
\newlabel{more-morphisms-lemma-flat-and-free-at-point-fibre@cref}{{[originallemma][0][]080Q}{[1][0][]source}{}{}{}}
\makeatother
\crefname{originallemma}{original lemma}{original lemmas}
\Crefname{originallemma}{Original lemma}{Original lemmas}
\makeatletter
\newlabel{morphisms-lemma-base-change-proper}{{01W4}{source}{}{sourcepointer-6}{}}
\newlabel{morphisms-lemma-base-change-proper@cref}{{[originallemma][0][]01W4}{[1][0][]source}{}{}{}}
\makeatother
\crefname{originallemma}{original lemma}{original lemmas}
\Crefname{originallemma}{Original lemma}{Original lemmas}
\makeatletter
\newlabel{more-morphisms-lemma-thicken-property-morphisms-cartesian}{{09ZW}{source}{}{sourcepointer-7}{}}
\newlabel{more-morphisms-lemma-thicken-property-morphisms-cartesian@cref}{{[originallemma][0][]09ZW}{[1][0][]source}{}{}{}}
\makeatother
\crefname{originallemma}{original lemma}{original lemmas}
\Crefname{originallemma}{Original lemma}{Original lemmas}
\makeatletter
\newlabel{morphisms-lemma-ample-base-change}{{0893}{source}{}{sourcepointer-8}{}}
\newlabel{morphisms-lemma-ample-base-change@cref}{{[originallemma][0][]0893}{[1][0][]source}{}{}{}}
\makeatother
\crefname{originallemma}{original lemma}{original lemmas}
\Crefname{originallemma}{Original lemma}{Original lemmas}
\makeatletter
\newlabel{more-morphisms-lemma-thicken-property-relatively-ample}{{0D2R}{source}{}{sourcepointer-9}{}}
\newlabel{more-morphisms-lemma-thicken-property-relatively-ample@cref}{{[originallemma][0][]0D2R}{[1][0][]source}{}{}{}}
\makeatother
\title{RS* descent for the polarized stack along pushouts along thickenings}
\author{The Stacks Project Authors}
\date{stacks-project:quot.tex:local-snapshot-20261012; GNU FDL 1.2 or later}
\begin{document}
\maketitle
\setcounter{subsection}{1}
\noindent\textit{Editorial note.} All delivered mathematical text is the author's own, copied from the pinned Stacks Project chapter file quot.tex; no statement or proof step was rewritten, reordered or removed. The counted span is the author environment and its own proof environment. References inside the span are handed to the assembler as navigation stubs that carry the canonical tag of the cited tag only; whether each of them is load-bearing is recorded on the stub itself.\par\medskip
\noindent\textit{Allowed standard prerequisites (editorial).}\begin{itemize}
\item Situation situation-coherent (canonical tag 08KB). cited at statement line 4320. Standing situation.
\item Lemma more-morphisms-lemma-pushout-along-thickening (canonical tag 07RT). cited at statement lines 4328-4329. The pushout square exists.
\item Lemma more-morphisms-lemma-equivalence-categories-schemes-over-pushout-flat (canonical tag 07RX). cited at proof lines 4340-4341. Flat-lfp schemes descend along the pushout.
\item Lemma more-morphisms-lemma-scheme-over-pushout-flat-modules (canonical tag 08KU). cited at proof lines 4353-4354. QCoh-flat modules descend, giving the Pic equivalence.
\item Lemma more-morphisms-lemma-flat-and-free-at-point-fibre (canonical tag 080Q). cited at proof lines 4380-4381. Invertibility checked on fibres.
\item Lemma morphisms-lemma-base-change-proper (canonical tag 01W4). cited at proof lines 4392-4393. Properness to the pieces.
\item Lemma more-morphisms-lemma-thicken-property-morphisms-cartesian (canonical tag 09ZW). cited at proof lines 4396-4397. Properness back to X'.
\item Lemma morphisms-lemma-ample-base-change (canonical tag 0893). cited at proof lines 4400-4401. Ampleness to the restrictions.
\item Lemma more-morphisms-lemma-thicken-property-relatively-ample (canonical tag 0D2R). cited at proof lines 4405-4406. Ampleness back to L'.
\end{itemize}\medskip
\ifdefined\hypertarget\hypertarget{sourcepointer-1}{}\fi
\noindent\textit{Source reference.} Original definition 08KB: sources/stacks-chow/tags.txt, source lines 10879--10879. Cited at statement line 4320: the lemma is phrased in this situation; load-bearing.\par\medskip
\ifdefined\hypertarget\hypertarget{sourcepointer-2}{}\fi
\noindent\textit{Source reference.} Original lemma 07RT: sources/stacks-chow/tags.txt, source lines 9878--9878. Cited at statement lines 4328-4329: existence of the pushout square; load-bearing.\par\medskip
\ifdefined\hypertarget\hypertarget{sourcepointer-3}{}\fi
\noindent\textit{Source reference.} Original lemma 07RX: sources/stacks-chow/tags.txt, source lines 9882--9882. Cited at proof lines 4340-4341 (flat-lfp equivalence): descent of the scheme X'/S' to the triple; load-bearing.\par\medskip
\ifdefined\hypertarget\hypertarget{sourcepointer-4}{}\fi
\noindent\textit{Source reference.} Original lemma 08KU: sources/stacks-chow/tags.txt, source lines 10897--10897. Cited at proof lines 4353-4354: QCoh-flat equivalence over the pushout, giving the Pic equivalence; load-bearing.\par\medskip
\ifdefined\hypertarget\hypertarget{sourcepointer-5}{}\fi
\noindent\textit{Source reference.} Original lemma 080Q: sources/stacks-chow/tags.txt, source lines 10190--10190. Cited at proof lines 4380-4381: F' finitely presented, invertibility checked on fibres; load-bearing for the small point.\par\medskip
\ifdefined\hypertarget\hypertarget{sourcepointer-6}{}\fi
\noindent\textit{Source reference.} Original lemma 01W4: sources/stacks-chow/tags.txt, source lines 2469--2469. Cited at proof lines 4392-4393: X' proper gives X and Y' proper; load-bearing.\par\medskip
\ifdefined\hypertarget\hypertarget{sourcepointer-7}{}\fi
\noindent\textit{Source reference.} Original lemma 09ZW: sources/stacks-chow/tags.txt, source lines 12614--12614. Cited at proof lines 4396-4397: X, Y' proper give X' proper; load-bearing.\par\medskip
\ifdefined\hypertarget\hypertarget{sourcepointer-8}{}\fi
\noindent\textit{Source reference.} Original lemma 0893: sources/stacks-chow/tags.txt, source lines 10483--10483. Cited at proof lines 4400-4401: L' ample gives restrictions ample; load-bearing.\par\medskip
\ifdefined\hypertarget\hypertarget{sourcepointer-9}{}\fi
\noindent\textit{Source reference.} Original lemma 0D2R: sources/stacks-chow/tags.txt, source lines 16391--16391. Cited at proof lines 4405-4406: restrictions ample give L' ample; load-bearing.\par\medskip
\begin{lemma}
\label{lemma-polarized-RS-star}
In Situation \ref{situation-coherent}. Let
$$
\xymatrix{
T \ar[r] \ar[d] & T' \ar[d] \\
S \ar[r] & S'
}
$$
be a pushout in the category of schemes where
$T \to T'$ is a thickening and $T \to S$ is affine, see
More on Morphisms, Lemma \ref{more-morphisms-lemma-pushout-along-thickening}.
Then the functor on fibre categories
$$
\Polarizedstack_{S'}
\longrightarrow
\Polarizedstack_S \times_{\Polarizedstack_T} \Polarizedstack_{T'}
$$
is an equivalence.
\end{lemma}

\begin{proof}
By More on Morphisms, Lemma
\ref{more-morphisms-lemma-equivalence-categories-schemes-over-pushout-flat}
there is an equivalence
$$
\textit{flat-lfp}_{S'}
\longrightarrow
\textit{flat-lfp}_S \times_{\textit{flat-lfp}_T} \textit{flat-lfp}_{T'}
$$
where $\textit{flat-lfp}_S$ signifies the category of schemes flat
and locally of finite presentation over $S$.
Let $X'/S'$ on the left hand side correspond to the triple
$(X/S, Y'/T', \varphi)$ on the right hand side.
Set $Y = T \times_{T'} Y'$ which is isomorphic with
$T \times_S X$ via $\varphi$. Then More on Morphisms, Lemma
\ref{more-morphisms-lemma-scheme-over-pushout-flat-modules}
shows that we have an equivalence
$$
\textit{QCoh-flat}_{X'/S'}
\longrightarrow
\textit{QCoh-flat}_{X/S}
\times_{\textit{QCoh-flat}_{Y/T}} \textit{QCoh-flat}_{Y'/T'}
$$
where $\textit{QCoh-flat}_{X/S}$ signifies the category of
quasi-coherent $\mathcal{O}_X$-modules flat over $S$.
Since $X \to S$, $Y \to T$, $X' \to S'$, $Y' \to T'$ are
flat, this will in particular apply to invertible modules
to give an equivalence of categories
$$
\textit{Pic}(X')
\longrightarrow
\textit{Pic}(X) \times_{\textit{Pic}(Y)} \textit{Pic}(Y')
$$
where $\textit{Pic}(X)$ signifies the category of invertible
$\mathcal{O}_X$-modules. There is a small point here:
one has to show that if an object $\mathcal{F}'$
of $\textit{QCoh-flat}_{X'/S'}$
pulls back to invertible modules on $X$ and $Y'$, then
$\mathcal{F}'$ is an invertible $\mathcal{O}_{X'}$-module.
It follows from the cited lemma that $\mathcal{F}'$
is an $\mathcal{O}_{X'}$-module of finite presentation.
By More on Morphisms, Lemma
\ref{more-morphisms-lemma-flat-and-free-at-point-fibre}
it suffices to check the restriction of
$\mathcal{F}'$ to fibres of $X' \to S'$ is invertible.
But the fibres of $X' \to S'$ are the same as the fibres
of $X \to S$ and hence these restrictions are invertible.

\medskip\noindent
Having said the above we obtain an equivalence of categories if we drop
the assumption (for the category of objects over $S$) that $X \to S$ be proper
and the assumption that $\mathcal{L}$ be ample.
Now it is clear that if $X' \to S'$ is proper, then
$X \to S$ and $Y' \to T'$ are proper (Morphisms, Lemma
\ref{morphisms-lemma-base-change-proper}).
Conversely, if $X \to S$ and $Y' \to T'$ are proper, then
$X' \to S'$ is proper by
More on Morphisms, Lemma
\ref{more-morphisms-lemma-thicken-property-morphisms-cartesian}.
Similarly, if $\mathcal{L}'$ is ample on $X'/S'$, then
$\mathcal{L}'|_X$ is ample on $X/S$ and
$\mathcal{L}'|_{Y'}$ is ample on $Y'/T'$
(Morphisms, Lemma \ref{morphisms-lemma-ample-base-change}).
Finally, if $\mathcal{L}'|_X$ is ample on $X/S$ and
$\mathcal{L}'|_{Y'}$ is ample on $Y'/T'$, then
$\mathcal{L}'$ is ample on $X'/S'$ by
More on Morphisms, Lemma
\ref{more-morphisms-lemma-thicken-property-relatively-ample}.
\end{proof}

\end{document}
